% Generated by Listas/make_statements.py from pratica-lecture-06.md, item 13. Do not edit.
\textbf{Source.} Muñoz-Garcia \& Toro-Gonzalez (2019), \emph{Strategy and Game Theory: Practice Exercises with Answers}, Ch. 6, Exercise 3 ``Temporary Punishments from Deviation'', pp.~282--283 {[}PDF 293--294{]}. Reproduced verbatim (statement only; the book's worked answer follows it).



\begin{quote}

Consider two firms competing a la Cournot, and facing linear inverse demand \(p(Q) = 100 - Q\), where \(Q = q_1 + q_2\) denotes aggregate output. For simplicity, assume that firms face a common marginal cost of production \(c = 10\).



\begin{enumerate}

\def\labelenumi{(\alph{enumi})}

\item

  \emph{Unrepeated game.} Find the equilibrium output each firm produces when competing a la Cournot (that is, when they simultaneously and independently choose their output levels) in the unrepeated version of the game (that is, when firms interact only once). In addition, find the profits that each firm earns in equilibrium.

\item

  \emph{Repeated game---Collusion.} Assume now that the CEOs from both companies meet to discuss a collusive agreement that would increase their profits. Set up the maximization problem that firms solve when maximizing their \emph{joint} profits (that is, the sum of profits for both firms). Find the output level that each firm should select to maximize joint profits. In addition, find the profits that each firm obtains in this collusive agreement.

\item

  \emph{Repeated game---Permanent punishment.} Consider a grim-trigger strategy in which every firm starts colluding in period 1, and it keeps doing so as long as both firms colluded in the past. Otherwise, every firm deviates to the Cournot equilibrium thereafter (that is, every firm produces the Nash equilibrium of the unrepeated game found in part a forever). In words, this says that the punishment of deviating from the collusive agreement is \emph{permanent}, since firms never return to the collusive outcome. For which discount factors this grim-trigger strategy can be sustained as the SPNE of the infinitely-repeated game?

\item

  \emph{Repeated game---Temporary punishment.} Consider now a ``modified'' grim-trigger strategy. Like in the grim-trigger strategy of part (c), every firm starts colluding in period 1, and it keeps doing so as long as both firms colluded in the past. However, if a deviation is detected by either firm, every firm deviates to the Cournot equilibrium during only 1 period, and then every firm returns to cooperation (producing the collusive output). Intuitively, this implies that the punishment of deviating from the collusive agreement is now \emph{temporary} (rather than permanent) since it lasts only one period. For which discount factors this ``modified'' grim-trigger strategy can be sustained as the SPNE of the infinitely-repeated game?

\item

  Consider again the temporary punishment in part (d), but assume now that it lasts for two periods. How are your results from part (d) affected? Interpret.

\item

  Consider again the temporary punishment in part (d), but assume now that it lasts for three periods. How are your results from part (d) affected? Interpret.

\end{enumerate}

\end{quote}

